% Generated draft: requires mathematical and source review.
\documentclass[11pt]{article}
\usepackage[margin=25mm]{geometry}
\usepackage{fontspec}
\setmainfont{Latin Modern Roman}
\usepackage{amsmath,amssymb,mathtools,unicode-math}
\setmathfont{Latin Modern Math}
\usepackage{enumitem,needspace,xurl,hyperref}
\setlength{\parindent}{0pt}
\setlength{\parskip}{5pt}
\setlength{\emergencystretch}{3em}
\newcommand{\N}{\mathbb{N}}
\newcommand{\Z}{\mathbb{Z}}
\newcommand{\Q}{\mathbb{Q}}
\newcommand{\R}{\mathbb{R}}
\newcommand{\C}{\mathbb{C}}
\newcommand{\F}{\mathbb{F}}
\newcommand{\A}{\mathbb{A}}
\newcommand{\PP}{\mathbb{P}}
\newcommand{\E}{\mathbb{E}}
\newcommand{\eps}{\varepsilon}
\newcommand{\dd}{\,\mathrm{d}}
\begin{document}
\section*{UnsolvedMath ID: 20003354}
\textbf{Source problem code:} catalogue entry
\subsection{Definitions and mathematical statement}
The literal record says only ``Conjecture 52 was established in [42].'' The aid explains that the omitted preceding qualification is ``for the above example,'' referring to the level-one modular curve. It does not establish the unreproduced general Conjecture 52. The following is the supplied developed statement for that example.

Let \(Y(1)=\operatorname{PSL}_2(\Z)\backslash\mathbb H\), with hyperbolic probability
\[
 m=\frac3\pi\frac{\dd x\,\dd y}{y^2}.
\]
For every negative imaginary-quadratic order discriminant \(D=f^2D_K\), where \(D_K\) is fundamental and \(f\geq1\), let \(\mu_D\) be the uniform probability on the full proper class-group/CM packet of that order. The complete packet, not an arbitrary chosen subset, is essential. The aid attributes equidistribution for fundamental discriminants to Duke's 1988 theorem [42], and extension to arbitrary imaginary-quadratic orders to Clozel--Ullmo:
\[
 \mu_D\Longrightarrow m\qquad(|D|\to\infty).
\]
The precise theorem hypotheses and references require checking.

For each \(n\), choose finitely supported weights \(w_{n,D}\geq0\), \(\sum_Dw_{n,D}=1\), and set
\[
 \nu_n=\sum_Dw_{n,D}\mu_D,\qquad
 W_n(B)=\sum_{|D|\leq B}w_{n,D}.
\]
The aid's exact stability criterion is
\[
 \nu_n\Longrightarrow m
 \quad\Longleftrightarrow\quad
 W_n(B)\longrightarrow0\text{ for every fixed finite }B.
\]
This allows arbitrary changing finite packet supports and nonuniform convex weights. It does not require every support discriminant to increase, only vanishing total weight in every bounded range. Conversely, positivity prevents a persistent bounded packet contribution from canceling, since bounded discriminant ranges give finitely many CM atoms while \(m\) has no atoms.

For every bounded continuous test \(a\), the reported two-term cutoff estimate may be written
\[
 |\nu_n(a)-m(a)|
 \leq2\|a\|_\infty W_n(B)
   +\sup_{|D|>B}|\mu_D(a)-m(a)|.
\]
Thus any discriminant windows with lower endpoints tending to infinity, and any normalized nonnegative weights supported in those windows, qualify. General windows qualify precisely by the stated weight-escape criterion; no signed-weight or arbitrary subpacket version is asserted.

For a fixed embedded standard cusp neighborhood \(C_Y\) of height \(Y\geq1\), the source's no-escape consequence is
\[
 \nu_n(C_Y)\longrightarrow m(C_Y)=\frac3{\pi Y}.
\]
The cusp boundary has zero hyperbolic measure. Sending \(Y\to\infty\) after this fixed-height limit gives the corresponding no-escape conclusion. The formula is not applied without a height convention or to an arbitrary subunit-height region. These are full-packet mixture consequences of the supplied modular-curve theorem, not a proof of general Conjecture 52.
\subsection{Short English statement}
The supplied modular-curve result concerns full CM packets. Their convex mixtures equidistribute exactly when bounded discriminant ranges receive vanishing total weight. The cutoff estimate also gives the stated fixed-cusp limit and no escape, without asserting the missing general conjecture.
\subsection{Sources}
\begin{itemize}
\item Local UnsolvedMath registry, permanent ID \(20003354\), exact source code \texttt{catalogue entry}; canonical \texttt{aim-other-notes.json}, \texttt{notes[160]}, AIM domain Other, registry category Miscellaneous, tag \texttt{conjecture}.
\item AIM workshop ``Emerging applications of measure rigidity,'' assigned Conjecture 52 status fragment: \url{https://aimath.org/WWN/measrigid/measrigid.pdf}. The aid restricts the assertion to the preceding level-one modular-curve example; the complete general conjecture is absent. No PDF was fetched.
\item Source-local [42], identified by the aid with Duke's 1988 theorem, and the later Clozel--Ullmo extension to arbitrary imaginary-quadratic orders. No titles, identifiers, or complete authoritative bibliographic entries are supplied. The original literature field is empty.
\item Research attempt 1: result \texttt{partial\_result}, problem status at run \texttt{solved}, rigor label \texttt{proved}, novelty confidence \texttt{low}; classification \texttt{SOLVED-IN-LITERATURE}, registry status \texttt{solved}. Created and updated 14 August 2026. These are supplied assessments for the specified example and developed mixture theorem, not a blanket statement about Conjecture 52.
\item Attribution: Ulam AI / UnsolvedMath Contributors, CC BY 4.0. The machine-generated research aid is not peer reviewed and requires independent verification before citation.
\end{itemize}
\end{document}
